%% kapa: exam=1 round=29 year=2018 period=1 form=A subject=civil_law subject_ko=민법 seq=22 number=22
%% source: Q-Net 감정평가사 기출문제 (cst00309) artlSeq=5202808 file=extracted/1008/5202808/2018년 제29회 감정평가사 1차시험 기출문제/2018년 제29회 감정평가사 1차시험 기출문제/1교시 A형.pdf
%% conversion: kordoc 4.24.51 -> tools/gen1.py (auto); review=proofread-sonnet; flags=none
\begin{question}{22}
전세권에 관한 설명으로 옳지 않은 것은? (다툼이 있으면 판례에 따름)
\begin{choices}
  \item 건물전세권이 법정갱신된 경우, 전세권자는 등기 없이도 전세권설정자나 그 목적물을 취득한 제3자에 대하여 갱신된 권리를 주장할 수 있다.
  \item 토지전세권의 존속기간을 약정하지 아니한 경우 각 당사자는 언제든지 상대방에 대하여 전세권의 소멸을 통고할 수 있다.
  \item 토지전세권의 존속기간을 1년 미만으로 정한 때에는 이를 1년으로 한다.
  \item 전세권자가 그 목적물의 성질에 의하여 정하여진 용법으로 이를 사용, 수익하지 아니한 경우에는 전세권설정자는 전세권의 소멸을 청구할 수 있다.
  \item 전세권 존속 중에는 장래에 그 전세권이 소멸하는 경우에 전세금반환채권이 발생하는 것을 조건으로 그 장래의 조건부채권을 양도할 수 있다.
\end{choices}
\end{question}
\answer{3}
